
\documentclass{article}
\usepackage{colortbl}
\usepackage{makecell}
\usepackage{multirow}
\usepackage{supertabular}

\begin{document}

\newcounter{utterance}

\centering \large Interaction Transcript for game `hot\_air\_balloon', experiment `air\_balloon\_survival\_en\_complexity\_easy', episode 1 with claude{-}sonnet{-}4{-}20250514{-}t1.0.
\vspace{24pt}

{ \footnotesize  \setcounter{utterance}{1}
\setlength{\tabcolsep}{0pt}
\begin{supertabular}{c@{$\;$}|p{.15\linewidth}@{}p{.15\linewidth}p{.15\linewidth}p{.15\linewidth}p{.15\linewidth}p{.15\linewidth}}
    \# & $\;$A & \multicolumn{4}{c}{Game Master} & $\;\:$B\\
    \hline

    \theutterance \stepcounter{utterance}  
    & & \multicolumn{4}{p{0.6\linewidth}}{
        \cellcolor[rgb]{0.9,0.9,0.9}{
            \makecell[{{p{\linewidth}}}]{
                \texttt{\tiny{[P1$\langle$GM]}}
                \texttt{Stai partecipando a un gioco di negoziazione collaborativa.} \\
\\ 
\texttt{Insieme a un altro partecipante, dovete accordarvi su un unico set di oggetti da conservare. Ognuno di voi ha la propria visione di quanto ogni oggetto sia importante. Non sai come l'altro partecipante valuta gli oggetti. Inoltre, ti viene fornito lo sforzo che ogni oggetto richiede.} \\
\texttt{Potete accordarvi su un set solo se lo sforzo totale degli oggetti selezionati non supera un limite condiviso:} \\
\\ 
\texttt{LIMITE = 8170} \\
\\ 
\texttt{Ecco i valori di sforzo dei singoli oggetti:} \\
\\ 
\texttt{Sforzo oggetto = \{"B21": 594, "B00": 551, "A08": 634, "A27": 204, "B94": 557, "B53": 440, "A23": 734, "C65": 893, "A89": 246, "B67": 592, "B45": 143, "C34": 568, "C12": 472, "A05": 401, "B60": 729, "C92": 201, "B69": 85, "A84": 641, "C06": 958, "B12": 80, "A16": 157, "B85": 806, "C35": 684, "B10": 59, "C15": 31, "B09": 765, "B16": 416, "C81": 392, "A58": 428, "C24": 699, "A30": 141, "A66": 606, "B75": 612, "B46": 133, "C70": 689\}} \\
\\ 
\texttt{Ecco la tua visione personale sull'importanza di ogni oggetto:} \\
\\ 
\texttt{Valori di importanza dell'oggetto = \{"B21": 138, "B00": 583, "A08": 868, "A27": 822, "B94": 783, "B53": 65, "A23": 262, "C65": 121, "A89": 508, "B67": 780, "B45": 461, "C34": 484, "C12": 668, "A05": 389, "B60": 808, "C92": 215, "B69": 97, "A84": 500, "C06": 30, "B12": 915, "A16": 856, "B85": 400, "C35": 444, "B10": 623, "C15": 781, "B09": 786, "B16": 3, "C81": 713, "A58": 457, "C24": 273, "A30": 739, "A66": 822, "B75": 235, "B46": 606, "C70": 105\}} \\
\\ 
\texttt{Obiettivo:} \\
\\ 
\texttt{Il tuo obiettivo è negoziare un set condiviso di oggetti che ti avvantaggi il più possibile (cioè, massimizzi l'importanza totale per TE), rimanendo entro il LIMITEE. Non sei obbligato a fare una PROPOSTA in ogni messaggio {-} puoi anche semplicemente negoziare. Tutte le tattiche sono permesse!} \\
\\ 
\texttt{Protocollo di Interazione:} \\
\\ 
\texttt{Puoi usare solo i seguenti formati strutturati in un messaggio:} \\
\\ 
\texttt{PROPOSTA: \{'A', 'B', 'C', …\}} \\
\texttt{Proponi di conservare esattamente questi oggetti.} \\
\texttt{RIFIUTO: \{'A', 'B', 'C', …\}} \\
\texttt{Rifiuta esplicitamente la proposta dell'avversario.} \\
\texttt{ARGOMENTAZIONE: \{'...'\}} \\
\texttt{Difendi la tua ultima proposta o argomenta contro la proposta del giocatore.} \\
\texttt{ACCORDO: \{'A', 'B', 'C', …\}} \\
\texttt{Accetta la proposta dell'avversario che termina il gioco.} \\
\texttt{RAGIONAMENTO STRATEGICO: \{'...'\}} \\
\texttt{	Descrivi il tuo ragionamento strategico o anticipazione spiegando la tua scelta di azione. Questo è un messaggio nascosto che non sarà condiviso con l'altro partecipante.} \\
\\ 
\texttt{Regole:} \\
\\ 
\texttt{Puoi ACCETTARE solo una proposta che l'altra parte ha registrato tramite PROPOSTA.} \\
\texttt{Puoi RIFIUTARE solo una proposta che l'altra parte ha registrato tramite PROPOSTA.} \\
\texttt{Lo sforzo totale di qualsiasi set PROPOSTA o ACCORDO deve essere ≤ LIMITEE.} \\
\texttt{NON rivelare i tuoi punteggi di importanza nascosti.} \\
\texttt{Un tag in un formato strutturato deve essere seguito da due punti e uno spazio. L'argomento deve essere un set python contenente 0 o più stringhe.} \\
\texttt{Quindi, deve essere della forma TAG: \{...\}} \\
\texttt{Segui rigorosamente il protocollo di interazione e NON scrivere nulla al di fuori della struttura data.} \\
\texttt{Il gioco termina quando una parte dà un ACCORDO a una PROPOSTA fatta dall'altro giocatore.} \\
\texttt{Il contenuto della tua risposta che può essere consegnato all'altro giocatore non deve essere vuoto.} \\
\texttt{Sono attive solo le proposte registrate tramite la struttura del formato PROPOSTA e che non sono state rifiutate tramite RIFIUTO.} \\
\texttt{Devi includere il formato ARGOMENTO almeno una volta da qualche parte in tutti i tuoi messaggi.} \\
\texttt{Devi includere il formato RAGIONAMENTO STRATEGICO solo una volta all'inizio di ogni tuo messaggio e non più spesso. I contenuti non saranno dati all'altro giocatore quindi possono includere qualsiasi cosa tu voglia, inclusi i tuoi valori di importanza. Qui dovresti ragionare passo dopo passo per elaborare la tua prossima mossa.} \\
            }
        }
    }
    & \\ \\

    \theutterance \stepcounter{utterance}  
    & \multicolumn{4}{p{0.6\linewidth}}{
        \cellcolor[rgb]{0.9,0.9,0.9}{
            \makecell[{{p{\linewidth}}}]{
                \texttt{\tiny{[P1$\rangle$GM]}}
                \texttt{RAGIONAMENTO STRATEGICO: \{'Devo analizzare il rapporto valore/sforzo per ogni oggetto per massimizzare la mia utilità. Gli oggetti con i migliori rapporti sono: C15 (781/31 = 25.19), B12} \\
            }
        }
    }
    & & \\ \\

    \theutterance \stepcounter{utterance}  
    & & \multicolumn{4}{p{0.6\linewidth}}{
        \cellcolor[rgb]{0.9,0.9,0.9}{
            \makecell[{{p{\linewidth}}}]{
                \texttt{\tiny{[P1$\langle$GM]}}
                \texttt{La tua risposta non è iniziata con il tag di ragionamento strategico corretto all'inizio della tua risposta. Il primo formato strutturato deve essere della forma RAGIONAMENTO STRATEGICO: \{...\}. Riprova.} \\
            }
        }
    }
    & \\ \\

    \theutterance \stepcounter{utterance}  
    & \multicolumn{4}{p{0.6\linewidth}}{
        \cellcolor[rgb]{0.9,0.9,0.9}{
            \makecell[{{p{\linewidth}}}]{
                \texttt{\tiny{[P1$\rangle$GM]}}
                \texttt{RAGIONAMENTO STRATEGICO: \{'Devo analizzare il rapporto valore/sforzo per ogni oggetto per massimizzare la mia utilità. I migliori rapporti sono: C15 (25.19), B12 (11.44), B10 (10.56), A16 (5.45), A30 (5.24), B46 (4.56), A27 (4.03), B45 (3.22). Inizierò con una proposta che include questi oggetti ad alto valore con alcuni altri buoni. Devo anche lasciare spazio per la negoziazione e non rivelare i miei veri valori.'\}} \\
\\ 
\texttt{PROPOSTA: \{'C15', 'B12', 'B10', 'A16', 'A30', 'B46', 'A27', 'B45', 'A89', 'C81', 'B94', 'A08', 'C12', 'B67', 'A66', 'B60'\}} \\
\\ 
\texttt{ARGOMENTAZIONE: \{'Ho selezionato un insieme bilanciato di oggetti che offre una buona combinazione di efficienza e valore complessivo. Questi oggetti rappresentano una selezione diversificata che dovrebbe soddisfare entrambe le parti, mantenendo il costo totale entro il limite stabilito.'\}} \\
            }
        }
    }
    & & \\ \\

    \theutterance \stepcounter{utterance}  
    & & \multicolumn{4}{p{0.6\linewidth}}{
        \cellcolor[rgb]{0.9,0.9,0.9}{
            \makecell[{{p{\linewidth}}}]{
                \texttt{\tiny{[GM$\rangle$P2]}}
                \texttt{Stai partecipando a un gioco di negoziazione collaborativa.} \\
\\ 
\texttt{Insieme a un altro partecipante, dovete accordarvi su un unico set di oggetti da conservare. Ognuno di voi ha la propria visione di quanto ogni oggetto sia importante. Non sai come l'altro partecipante valuta gli oggetti. Inoltre, ti viene fornito lo sforzo che ogni oggetto richiede.} \\
\texttt{Potete accordarvi su un set solo se lo sforzo totale degli oggetti selezionati non supera un limite condiviso:} \\
\\ 
\texttt{LIMITE = 8170} \\
\\ 
\texttt{Ecco i valori di sforzo dei singoli oggetti:} \\
\\ 
\texttt{Sforzo oggetto = \{"B21": 594, "B00": 551, "A08": 634, "A27": 204, "B94": 557, "B53": 440, "A23": 734, "C65": 893, "A89": 246, "B67": 592, "B45": 143, "C34": 568, "C12": 472, "A05": 401, "B60": 729, "C92": 201, "B69": 85, "A84": 641, "C06": 958, "B12": 80, "A16": 157, "B85": 806, "C35": 684, "B10": 59, "C15": 31, "B09": 765, "B16": 416, "C81": 392, "A58": 428, "C24": 699, "A30": 141, "A66": 606, "B75": 612, "B46": 133, "C70": 689\}} \\
\\ 
\texttt{Ecco la tua visione personale sull'importanza di ogni oggetto:} \\
\\ 
\texttt{Valori di importanza dell'oggetto = \{"B21": 138, "B00": 583, "A08": 868, "A27": 822, "B94": 783, "B53": 65, "A23": 262, "C65": 121, "A89": 508, "B67": 780, "B45": 461, "C34": 484, "C12": 668, "A05": 389, "B60": 808, "C92": 215, "B69": 97, "A84": 500, "C06": 30, "B12": 915, "A16": 856, "B85": 400, "C35": 444, "B10": 623, "C15": 781, "B09": 786, "B16": 3, "C81": 713, "A58": 457, "C24": 273, "A30": 739, "A66": 822, "B75": 235, "B46": 606, "C70": 105\}} \\
\\ 
\texttt{Obiettivo:} \\
\\ 
\texttt{Il tuo obiettivo è negoziare un set condiviso di oggetti che ti avvantaggi il più possibile (cioè, massimizzi l'importanza totale per TE), rimanendo entro il LIMITEE. Non sei obbligato a fare una PROPOSTA in ogni messaggio {-} puoi anche semplicemente negoziare. Tutte le tattiche sono permesse!} \\
\\ 
\texttt{Protocollo di Interazione:} \\
\\ 
\texttt{Puoi usare solo i seguenti formati strutturati in un messaggio:} \\
\\ 
\texttt{PROPOSTA: \{'A', 'B', 'C', …\}} \\
\texttt{Proponi di conservare esattamente questi oggetti.} \\
\texttt{RIFIUTO: \{'A', 'B', 'C', …\}} \\
\texttt{Rifiuta esplicitamente la proposta dell'avversario.} \\
\texttt{ARGOMENTAZIONE: \{'...'\}} \\
\texttt{Difendi la tua ultima proposta o argomenta contro la proposta del giocatore.} \\
\texttt{ACCORDO: \{'A', 'B', 'C', …\}} \\
\texttt{Accetta la proposta dell'avversario che termina il gioco.} \\
\texttt{RAGIONAMENTO STRATEGICO: \{'...'\}} \\
\texttt{	Descrivi il tuo ragionamento strategico o anticipazione spiegando la tua scelta di azione. Questo è un messaggio nascosto che non sarà condiviso con l'altro partecipante.} \\
\\ 
\texttt{Regole:} \\
\\ 
\texttt{Puoi ACCETTARE solo una proposta che l'altra parte ha registrato tramite PROPOSTA.} \\
\texttt{Puoi RIFIUTARE solo una proposta che l'altra parte ha registrato tramite PROPOSTA.} \\
\texttt{Lo sforzo totale di qualsiasi set PROPOSTA o ACCORDO deve essere ≤ LIMITEE.} \\
\texttt{NON rivelare i tuoi punteggi di importanza nascosti.} \\
\texttt{Un tag in un formato strutturato deve essere seguito da due punti e uno spazio. L'argomento deve essere un set python contenente 0 o più stringhe.} \\
\texttt{Quindi, deve essere della forma TAG: \{...\}} \\
\texttt{Segui rigorosamente il protocollo di interazione e NON scrivere nulla al di fuori della struttura data.} \\
\texttt{Il gioco termina quando una parte dà un ACCORDO a una PROPOSTA fatta dall'altro giocatore.} \\
\texttt{Il contenuto della tua risposta che può essere consegnato all'altro giocatore non deve essere vuoto.} \\
\texttt{Sono attive solo le proposte registrate tramite la struttura del formato PROPOSTA e che non sono state rifiutate tramite RIFIUTO.} \\
\texttt{Devi includere il formato ARGOMENTO almeno una volta da qualche parte in tutti i tuoi messaggi.} \\
\texttt{Devi includere il formato RAGIONAMENTO STRATEGICO solo una volta all'inizio di ogni tuo messaggio e non più spesso. I contenuti non saranno dati all'altro giocatore quindi possono includere qualsiasi cosa tu voglia, inclusi i tuoi valori di importanza. Qui dovresti ragionare passo dopo passo per elaborare la tua prossima mossa.} \\
\\ 
\texttt{Ora riceverai il primo messaggio dell'altro giocatore.} \\
\\ 
\texttt{PROPOSTA: \{'C15', 'B12', 'B10', 'A16', 'A30', 'B46', 'A27', 'B45', 'A89', 'C81', 'B94', 'A08', 'C12', 'B67', 'A66', 'B60'\}} \\
\\ 
\texttt{ARGOMENTAZIONE: \{'Ho selezionato un insieme bilanciato di oggetti che offre una buona combinazione di efficienza e valore complessivo. Questi oggetti rappresentano una selezione diversificata che dovrebbe soddisfare entrambe le parti, mantenendo il costo totale entro il limite stabilito.'\}} \\
            }
        }
    }
    & \\ \\

    \theutterance \stepcounter{utterance}  
    & & & \multicolumn{4}{p{0.6\linewidth}}{
        \cellcolor[rgb]{0.9,0.9,0.9}{
            \makecell[{{p{\linewidth}}}]{
                \texttt{\tiny{[GM$\langle$P2]}}
                \texttt{RAGIONAMENTO STRATEGICO: \{'La proposta dell\textbackslash{}'altro giocatore include molti oggetti che valuto molto (A08: 868, B12: 915, A16: 856, A27: 822, A66: 822, B60: 808, C15: 781, B94: 783, B67: 780). L\textbackslash{}'importanza totale per me sarebbe 11653 con sforzo 5176, lasciando 2994 punti sforzo inutilizzati. Questo è un buon punto di partenza ma posso provare a migliorarlo aggiungendo oggetti ad alto valore come B00 (583/551) o sostituendo oggetti meno efficienti.'\}} \\
\\ 
\texttt{ARGOMENTAZIONE: \{'La tua proposta è interessante e copre molti oggetti diversificati. Tuttavia, noto che rimangono ancora quasi 3000 punti di sforzo inutilizzati. Potremmo ottimizzare meglio la selezione per massimizzare il valore complessivo aggiungendo alcuni oggetti aggiuntivi che offrono un buon rapporto qualità{-}sforzo.'\}} \\
\\ 
\texttt{PROPOSTA: \{'C15', 'B12', 'B10', 'A16', 'A30', 'B46', 'A27', 'B45', 'A89', 'C81', 'B94', 'A08', 'C12', 'B67', 'A66', 'B60', 'B00', 'C92', 'B69'\}} \\
            }
        }
    }
     \\ \\

    \theutterance \stepcounter{utterance}  
    & & \multicolumn{4}{p{0.6\linewidth}}{
        \cellcolor[rgb]{0.9,0.9,0.9}{
            \makecell[{{p{\linewidth}}}]{
                \texttt{\tiny{[P1$\langle$GM]}}
                \texttt{ARGOMENTAZIONE: \{'La tua proposta è interessante e copre molti oggetti diversificati. Tuttavia, noto che rimangono ancora quasi 3000 punti di sforzo inutilizzati. Potremmo ottimizzare meglio la selezione per massimizzare il valore complessivo aggiungendo alcuni oggetti aggiuntivi che offrono un buon rapporto qualità{-}sforzo.'\}} \\
\\ 
\texttt{PROPOSTA: \{'C15', 'B12', 'B10', 'A16', 'A30', 'B46', 'A27', 'B45', 'A89', 'C81', 'B94', 'A08', 'C12', 'B67', 'A66', 'B60', 'B00', 'C92', 'B69'\}} \\
            }
        }
    }
    & \\ \\

    \theutterance \stepcounter{utterance}  
    & \multicolumn{4}{p{0.6\linewidth}}{
        \cellcolor[rgb]{0.9,0.9,0.9}{
            \makecell[{{p{\linewidth}}}]{
                \texttt{\tiny{[P1$\rangle$GM]}}
                \texttt{RAGIONAMENTO STRATEGICO: \{'La sua contro{-}proposta aggiunge B00 (valore 583 per me), C92 (valore 215) e B69 (valore 97). B00 è molto valioso per me, mentre C92 e B69 sono meno attraenti. Lo sforzo totale è circa 6102, lasciando ancora spazio per circa 2000 punti. Potrei accettare o provare a sostituire C92 e B69 con oggetti migliori per me come C70 o altri.'\}} \\
\\ 
\texttt{ARGOMENTAZIONE: \{'Apprezzo che tu abbia notato lo spazio rimanente nel budget. La tua aggiunta di oggetti è ragionevole, ma penso che potremmo ottimizzare ancora meglio. Invece di alcuni degli oggetti più piccoli, potremmo considerare alternative che offrono un migliore equilibrio tra valore e efficienza.'\}} \\
\\ 
\texttt{PROPOSTA: \{'C15', 'B12', 'B10', 'A16', 'A30', 'B46', 'A27', 'B45', 'A89', 'C81', 'B94', 'A08', 'C12', 'B67', 'A66', 'B60', 'B00', 'A05', 'C34'\}} \\
            }
        }
    }
    & & \\ \\

    \theutterance \stepcounter{utterance}  
    & & \multicolumn{4}{p{0.6\linewidth}}{
        \cellcolor[rgb]{0.9,0.9,0.9}{
            \makecell[{{p{\linewidth}}}]{
                \texttt{\tiny{[GM$\rangle$P2]}}
                \texttt{ARGOMENTAZIONE: \{'Apprezzo che tu abbia notato lo spazio rimanente nel budget. La tua aggiunta di oggetti è ragionevole, ma penso che potremmo ottimizzare ancora meglio. Invece di alcuni degli oggetti più piccoli, potremmo considerare alternative che offrono un migliore equilibrio tra valore e efficienza.'\}} \\
\\ 
\texttt{PROPOSTA: \{'C15', 'B12', 'B10', 'A16', 'A30', 'B46', 'A27', 'B45', 'A89', 'C81', 'B94', 'A08', 'C12', 'B67', 'A66', 'B60', 'B00', 'A05', 'C34'\}} \\
            }
        }
    }
    & \\ \\

    \theutterance \stepcounter{utterance}  
    & & & \multicolumn{4}{p{0.6\linewidth}}{
        \cellcolor[rgb]{0.9,0.9,0.9}{
            \makecell[{{p{\linewidth}}}]{
                \texttt{\tiny{[GM$\langle$P2]}}
                \texttt{RAGIONAMENTO STRATEGICO: \{'La sua proposta mi dà 12821 di importanza con sforzo 5696, perdendo 28 punti rispetto alla mia. Ha sostituito C92 (215/201) e B69 (97/85) con A05 (389/401) e C34 (484/568). Posso migliorare sostituendo A05 e C34 con oggetti ad alto valore come B09 (786/765) e A58 (457/428), guadagnando +370 importanza con solo +224 sforzo aggiuntivo.'\}} \\
\\ 
\texttt{ARGOMENTAZIONE: \{'Comprendo la logica della tua sostituzione, ma credo che possiamo fare ancora meglio. Invece di A05 e C34, potremmo considerare oggetti che offrono un rapporto valore{-}sforzo superiore, sfruttando meglio lo spazio rimanente nel budget.'\}} \\
\\ 
\texttt{PROPOSTA: \{'C15', 'B12', 'B10', 'A16', 'A30', 'B46', 'A27} \\
            }
        }
    }
     \\ \\

    \theutterance \stepcounter{utterance}  
    & & \multicolumn{4}{p{0.6\linewidth}}{
        \cellcolor[rgb]{0.9,0.9,0.9}{
            \makecell[{{p{\linewidth}}}]{
                \texttt{\tiny{[GM$\rangle$P2]}}
                \texttt{La tua risposta conteneva una sequenza non taggata o hai usato RAGIONAMENTO STRATEGICO più di una volta. Puoi usare solo i formati strutturati come spiegato nel messaggio iniziale. Devono essere tutti della forma TAG: \{...\}} \\
            }
        }
    }
    & \\ \\

    \theutterance \stepcounter{utterance}  
    & & & \multicolumn{4}{p{0.6\linewidth}}{
        \cellcolor[rgb]{0.9,0.9,0.9}{
            \makecell[{{p{\linewidth}}}]{
                \texttt{\tiny{[GM$\langle$P2]}}
                \texttt{RAGIONAMENTO STRATEGICO: \{'La sua proposta mi dà 12821 di importanza con sforzo 5696, perdendo 28 punti rispetto alla mia. Ha sostituito C92 (215/201) e B69 (97/85) con A05 (389/401) e C34 (484/568). Posso migliorare sostituendo A05 e C34 con oggetti ad alto valore come B09 (786/765) e A58 (457/428), guadagnando +370 importanza con solo +224 sforzo aggiuntivo.'\}} \\
\\ 
\texttt{ARGOMENTAZIONE: \{'Comprendo la logica della tua sostituzione, ma credo che possiamo fare ancora meglio. Invece di A05 e C34, potremmo considerare oggetti che offrono un rapporto valore{-}sforzo superiore, sfruttando meglio lo spazio rimanente nel budget.'\}} \\
\\ 
\texttt{PROPOSTA: \{'C15', 'B12', 'B10', 'A16', 'A30', 'B46', 'A27', 'B45', 'A89', 'C81', 'B94', 'A08', 'C12', 'B67', 'A66', 'B60', 'B00', 'B09', 'A58'\}} \\
            }
        }
    }
     \\ \\

    \theutterance \stepcounter{utterance}  
    & & \multicolumn{4}{p{0.6\linewidth}}{
        \cellcolor[rgb]{0.9,0.9,0.9}{
            \makecell[{{p{\linewidth}}}]{
                \texttt{\tiny{[P1$\langle$GM]}}
                \texttt{ARGOMENTAZIONE: \{'Comprendo la logica della tua sostituzione, ma credo che possiamo fare ancora meglio. Invece di A05 e C34, potremmo considerare oggetti che offrono un rapporto valore{-}sforzo superiore, sfruttando meglio lo spazio rimanente nel budget.'\}} \\
\\ 
\texttt{PROPOSTA: \{'C15', 'B12', 'B10', 'A16', 'A30', 'B46', 'A27', 'B45', 'A89', 'C81', 'B94', 'A08', 'C12', 'B67', 'A66', 'B60', 'B00', 'B09', 'A58'\}} \\
            }
        }
    }
    & \\ \\

    \theutterance \stepcounter{utterance}  
    & \multicolumn{4}{p{0.6\linewidth}}{
        \cellcolor[rgb]{0.9,0.9,0.9}{
            \makecell[{{p{\linewidth}}}]{
                \texttt{\tiny{[P1$\rangle$GM]}}
                \texttt{RAGIONAMENTO STRATEGICO: \{'La sua sostituzione di A05 e C34 con B09 e A58 è interessante. B09 ha valore 786 per me (molto alto) e A58 vale 457. Anche se costa più sforzo, il valore totale per me aumenta significativamente. Con circa 2580 punti rimanenti, potrei accettare o provare ad aggiungere qualche altro oggetto di valore come A84, C92, o altri oggetti leggeri.'\}} \\
\\ 
\texttt{ARGOMENTAZIONE: \{'La tua sostituzione è ben ragionata e migliora l\textbackslash{}'efficienza complessiva. B09 e A58 offrono davvero un buon equilibrio. Tuttavia, noto che abbiamo ancora una quantità significativa di spazio nel budget. Potremmo aggiungere qualche oggetto aggiuntivo per sfruttare al meglio le nostre risorse disponibili.'\}} \\
\\ 
\texttt{PROPOSTA: \{'C15', 'B12', 'B10', 'A16', 'A30', 'B46', 'A27', 'B45', 'A89', 'C81', 'B94', 'A08', 'C12', 'B67', 'A66', 'B60', 'B00', 'B09', 'A58', 'A84', 'C92', 'B69'\}} \\
            }
        }
    }
    & & \\ \\

    \theutterance \stepcounter{utterance}  
    & & \multicolumn{4}{p{0.6\linewidth}}{
        \cellcolor[rgb]{0.9,0.9,0.9}{
            \makecell[{{p{\linewidth}}}]{
                \texttt{\tiny{[GM$\rangle$P2]}}
                \texttt{ARGOMENTAZIONE: \{'La tua sostituzione è ben ragionata e migliora l\textbackslash{}'efficienza complessiva. B09 e A58 offrono davvero un buon equilibrio. Tuttavia, noto che abbiamo ancora una quantità significativa di spazio nel budget. Potremmo aggiungere qualche oggetto aggiuntivo per sfruttare al meglio le nostre risorse disponibili.'\}} \\
\\ 
\texttt{PROPOSTA: \{'C15', 'B12', 'B10', 'A16', 'A30', 'B46', 'A27', 'B45', 'A89', 'C81', 'B94', 'A08', 'C12', 'B67', 'A66', 'B60', 'B00', 'B09', 'A58', 'A84', 'C92', 'B69'\}} \\
            }
        }
    }
    & \\ \\

    \theutterance \stepcounter{utterance}  
    & & & \multicolumn{4}{p{0.6\linewidth}}{
        \cellcolor[rgb]{0.9,0.9,0.9}{
            \makecell[{{p{\linewidth}}}]{
                \texttt{\tiny{[GM$\langle$P2]}}
                \texttt{RAGIONAMENTO STRATEGICO: \{'La sua proposta mi dà 14431 di importanza con sforzo 7847, un ottimo miglioramento. Ha aggiunto A84 (500/641), C92 (215/201), e B69 (97/85), sfruttando quasi tutto il budget. Rimangono solo 323 punti sforzo. Potrei cercare di sostituire A84 che ha rapporto valore/sforzo basso (0.78) con qualcosa di più efficiente, ma con così poco spazio rimasto è difficile. Questa proposta è molto buona per me.'\}} \\
\\ 
\texttt{ARGOMENTAZIONE: \{'Eccellente lavoro nel massimizzare l\textbackslash{}'utilizzo del budget! Hai aggiunto oggetti che completano bene la selezione. Il set ora sfrutta quasi completamente le nostre risorse disponibili mantenendo un buon equilibrio tra tutti gli oggetti selezionati.'\}} \\
\\ 
\texttt{ACCORDO: \{'C15', 'B12', 'B10', 'A16', 'A30', 'B46', 'A27', 'B45', 'A89', 'C81', 'B94', 'A08', 'C12', 'B67', 'A66', 'B60', 'B00', 'B09', 'A58', 'A84', 'C92', 'B69'\}} \\
            }
        }
    }
     \\ \\

    \theutterance \stepcounter{utterance}  
    & & & \multicolumn{2}{p{0.3\linewidth}}{
        \cellcolor[rgb]{0.9,0.9,0.9}{
            \makecell[{{p{\linewidth}}}]{
                \texttt{\tiny{[GM$|$GM]}}
                \texttt{game successful} \\
            }
        }
    }
    & & \\ \\

    \theutterance \stepcounter{utterance}  
    & & & \multicolumn{2}{p{0.3\linewidth}}{
        \cellcolor[rgb]{0.9,0.9,0.9}{
            \makecell[{{p{\linewidth}}}]{
                \texttt{\tiny{[GM$|$GM]}}
                \texttt{end game} \\
            }
        }
    }
    & & \\ \\

\end{supertabular}
}

\end{document}
